% Relatório TP2 -- DCC207 Algoritmos 2 -- modelo SBC
\documentclass[12pt]{article}

\usepackage{sbc-template}
\usepackage{graphicx,url}
\usepackage[utf8]{inputenc}
\usepackage[T1]{fontenc}
\usepackage[brazil]{babel}
\usepackage{amsmath,amssymb}
\usepackage{booktabs}
\usepackage{xcolor}

% Resultados gerados por `python -m src.analysis`
\newcommand{\resdir}{../results/analysis}
\graphicspath{{\resdir/figs/}{figs/}}
\newcommand{\restable}[1]{\IfFileExists{\resdir/tables/#1.tex}{\resizebox{\linewidth}{!}{\input{\resdir/tables/#1.tex}}}{\textcolor{red}{[tabela #1 ainda não gerada]}}}
\newcommand{\todo}[1]{\textcolor{red}{[TODO: #1]}}

\sloppy

\title{Branch-and-Bound, Twice-around-the-tree e Christofides para o TSP Euclidiano:\\
análise empírica da poda, complexidade e da sensibilidade}

\author{Tomás Almeida Moreira}

\address{Departamento de Ciência da Computação\\
  Universidade Federal de Minas Gerais (UFMG)\\
  Belo Horizonte -- MG -- Brasil
  \email{tomasmoreira@ufmg.br}
}

\begin{document}

\maketitle

\begin{abstract}
  This work empirically evaluates an exact Branch-and-Bound and two approximation
  algorithms (Twice-around-the-tree and Christofides) for the Euclidean Traveling
  Salesman Problem on 15 TSPLIB-derived instances with 12 to 70 vertices. With a
  Held-Karp lower bound, the Branch-and-Bound solved all 75 runs within 30 minutes,
  whereas an MST bound timed out on 5 instances, from $N=22$ on. Among five
  upper-bound warm-starts, a convex-hull insertion tour reduced expanded nodes 12-fold
  and total time 2.1-fold; nearest neighbor and Twice-around-the-tree gave no gain,
  because the depth-first search finds an equivalent tour on its first descent.
  Profiling shows that the minimum-weight matching accounts for over 94\% of
  Christofides' time, with empirical growth $N^{2.7}$. Geometric features correlate
  with the approximation gap, but do not predict it out of sample.
\end{abstract}

\begin{resumo}
  Este trabalho avalia empiricamente um algoritmo exato (Branch-and-Bound) e dois
  algoritmos aproximados (Twice-around-the-tree e Christofides) para o Problema do
  Caixeiro Viajante euclidiano, em 15 instâncias derivadas da TSPLIB com 12 a 70
  vértices. Medimos o efeito de cinco estratégias de inicialização do limite superior
  (\textit{warm-start}) e de três cotas inferiores sobre a árvore de busca, ajustamos
  curvas de complexidade empírica às sub-rotinas instrumentadas e relacionamos
  características geométricas das instâncias ao gap de aproximação. Com a cota de
  Held-Karp, o B\&B resolveu as 75 execuções em menos de 30\,min, enquanto a cota de
  árvore geradora mínima esgotou o tempo em 5 instâncias, a partir de $N=22$. A
  inicialização por fecho convexo com inserção reduziu os nós expandidos em 12 vezes e o
  tempo total em 2,1 vezes; vizinho mais próximo e Twice-around-the-tree não trouxeram
  ganho, pois a busca em profundidade encontra um tour equivalente na primeira descida.
  O emparelhamento perfeito responde por mais de 94\% do tempo do Christofides, com
  crescimento empírico $N^{2,7}$. As \textit{features} geométricas se correlacionam com o
  gap, mas não o predizem fora da amostra.
\end{resumo}

% ============================================================================
\section{Introdução}

O Problema do Caixeiro Viajante (TSP) pede o ciclo hamiltoniano de menor custo em um
grafo completo ponderado. Sua versão de decisão é NP-completa, e mesmo o caso
euclidiano é NP-difícil~\cite{applegate2006}. Há dois caminhos usuais. Os algoritmos
exatos têm tempo exponencial no pior caso, mas o tempo real depende muito da qualidade
da poda. Os algoritmos aproximados rodam em tempo polinomial e têm garantia de
desempenho quando os custos obedecem à desigualdade triangular: fator 2 para o
\textit{Twice-around-the-tree} e 1{,}5 para o algoritmo de
Christofides~\cite{christofides1976}.

Além de implementar os algoritmos, medimos como as implementações se comportam: quanto
tempo e memória cada sub-rotina consome, quanto a escolha do limite superior inicial e
da cota inferior muda a árvore de busca, e que características geométricas das
instâncias pioram ou melhoram as aproximações. As contribuições são:
\begin{itemize}
  \item uma implementação de Branch-and-Bound em C++11 com três cotas inferiores
        (ingênua, árvore geradora mínima e Held-Karp) e contagem de nós
        expandidos e podados, inclusive em execuções interrompidas pelo limite de tempo;
  \item a comparação de cinco estratégias de \textit{warm-start} quanto à poda e ao
        custo total (heurística + busca);
  \item modelos de regressão para a complexidade empírica de cada sub-rotina e um
        modelo de sensibilidade do gap de aproximação às \textit{features} das instâncias.
\end{itemize}

A Seção~\ref{sec:impl} descreve a implementação e as decisões de auditoria; a
Seção~\ref{sec:metodo}, a metodologia experimental; a Seção~\ref{sec:res}, os resultados;
e a Seção~\ref{sec:concl} conclui.

% ============================================================================
\section{Implementação, Profiling e Auditoria}\label{sec:impl}

O núcleo exato foi escrito em C++11 e é invocado como subprocesso por um
\textit{harness} em Python~3, que contém as heurísticas, as aproximações, a extração
de \textit{features} e a análise estatística. O C++ ficou com a parte cara: o
Branch-and-Bound visita milhões de nós, cada um com um cálculo de cota de custo
$O(n^2)$, e em Python puro cada operação desse cálculo seria interpretada.
A IA generativa foi usada para gerar o esqueleto de parte do código, mas não para 
escrever a lógica dos algoritmos. O código final foi revisado e refatorado manualmente, 
e a auditoria documenta as decisões de projeto. Devido a uma constrição de tempo (eu não 
podia deixar meu computador rodando por vários dias, pois preciso dele para trabalhar), a 
refatoração do código para que ele pudesse rodar mais rápido usando o paralelismo, foi 
feita por IA, uma vez que eu não conseguiria, em tempo hábil, aprender esse conhecimento 
para implementar algo que eu nunca tinha feito antes (dividir processamento entre vários núcleos),
e entregar o trabalho no prazo. Para mais detalhes, ver a Tabela~\ref{tab:auditoria}. E o anexo de 
prompts.

\subsection{Branch-and-Bound}

\paragraph{Estratégia de busca.} Usamos busca em profundidade (DFS) com o vértice
inicial fixo em $0$ e filhos ordenados pela distância ao último vértice do caminho
(vizinho mais próximo primeiro). Descartamos a busca \textit{best-first} porque sua fila
de prioridade cresce exponencialmente e esgotaria a memória da máquina (7\,GB) por volta
de $N \approx 30$. A DFS guarda só o caminho corrente, em $O(n)$. O estado é um vetor de
caminho e um vetor de visitados, atualizados na descida e desfeitos no retorno da
recursão; nenhum nó copia matriz ou outra estrutura.

\paragraph{Cotas inferiores.} Para um caminho parcial $v_0 \to \dots \to v_k$ com
conjunto $U$ de vértices não visitados:
\begin{itemize}
  \item \textbf{naive}: $c(\text{caminho}) + \frac{1}{2}\big(m_1(v_k) + m_1(v_0) +
        \sum_{u \in U} (m_1(u) + m_2(u))\big)$, em que $m_1, m_2$ são as duas menores
        arestas incidentes~\cite{levitin2012}; custo $O(1)$ por nó com soma incremental;
  \item \textbf{mst}: o restante do ciclo é um caminho hamiltoniano $v_k \to U \to v_0$;
        removidas as arestas das pontas, sobra uma árvore geradora de $U$. Logo
        $LB = c(\text{caminho}) + \text{MST}(U) + \min_{u} d(v_k,u) + \min_{u} d(u,v_0)$,
        calculada com Prim denso $O(|U|^2)$;
  \item \textbf{hk}: a mesma cota sobre custos $d'(i,j) = d(i,j) + \pi_i + \pi_j$, com
        as penalidades $\pi$ de Held-Karp~\cite{heldkarp1970,heldkarp1971}, calculadas
        uma vez, na raiz, por subgradiente sobre 1-trees. A cota vale para
        qualquer $\pi$:
        $LB = LB'(\text{resto}) - \pi_{v_k} - \pi_{v_0} - 2\sum_{u\in U}\pi_u$.
\end{itemize}
O passo do subgradiente não depende do limite superior. Assim, o \textit{warm-start}
afeta a busca só pela poda, que é a variável que queremos medir. Como os custos são
inteiros, um nó é podado quando $\lceil LB \rceil \ge UB$.

\paragraph{Instrumentação.} O binário reporta nós expandidos, nós podados, folhas,
atualizações do limite superior, profundidade máxima, cota na raiz, tempo até a melhor
solução, tempo estimado na função de cota (amostragem de 1 a cada 256 chamadas) e pico
de memória (\texttt{getrusage}). O limite de 30\,min é verificado dentro da busca, de
modo que execuções interrompidas (registradas como NA) preservam os contadores.

A medida de memória saiu errada na rodada principal. No Linux, \texttt{ru\_maxrss}
sobrevive ao \texttt{exec}, e o binário lançado pelo \textit{harness} herda o pico do
processo Python que o criou: todas as 150 execuções reportaram de 35 a 36\,MB, sem
relação com $N$. Medido isoladamente, o binário ocupa 3,0\,MB numa instância com $N=3$
e 3,1\,MB na eil51 (44\,554 nós expandidos). Quase todo esse pico vem do
\textit{runtime} do C++, e o valor praticamente constante é o esperado para uma DFS com
memória $O(n)$.

\subsection{Heurísticas de inicialização e algoritmos aproximados}

\begin{itemize}
  \item \textbf{Trivial}: $0 \to 1 \to \dots \to n-1 \to 0$.
  \item \textbf{Vizinho mais próximo} a partir do vértice $0$, $O(n^2)$, com máscara
        vetorizada em NumPy.
  \item \textbf{Fecho convexo + inserção mais barata}: fecho por \textit{monotone
        chain}~\cite{andrew1979}, $O(n \log n)$; depois insere-se, a cada passo, o par
        (vértice, aresta) de menor acréscimo, $O(n^2)$ por passo e $O(n^3)$ no total.
  \item \textbf{Twice-around-the-tree}: MST por Prim denso~\cite{prim1957} em $O(n^2)$
        (melhor que Prim com heap, $O(n^2 \log n)$, pois o grafo é completo), seguida de
        pré-ordem a partir de $0$ (equivalente ao circuito euleriano com atalhos).
  \item \textbf{Christofides}: MST, vértices de grau ímpar, emparelhamento perfeito de
        peso mínimo (blossom de Edmonds~\cite{edmonds1965} via
        NetworkX~\cite{networkx2008}), circuito euleriano (Hierholzer iterativo) e atalhos.
\end{itemize}

\subsection{Profiling e auditoria}

Cada sub-rotina é cronometrada separadamente (\texttt{perf\_counter}, mediana de 5
repetições). O pico de memória é medido com \texttt{tracemalloc} numa passada
separada, porque o rastreamento de alocações deixa o código de 2 a 3 vezes mais lento.
No Christofides, a construção do grafo do emparelhamento e o emparelhamento somam 94\%
do tempo em $N=50$, 99\% em $N=200$ e praticamente 100\% em $N=1000$. Nesse último
tamanho, o Christofides leva 66\,s, e o Twice-around-the-tree, que usa a mesma MST,
leva 0,03\,s. As duas etapas também ocupam quase toda a memória, porque o NetworkX
representa o grafo completo dos $k$ vértices ímpares como $O(k^2)$ objetos Python: são
cerca de 39\,MB em $N=1000$, contra 68\,KB da MST.

\begin{table}[ht]
\centering
\caption{Decisões de auditoria e refatorações.}
\label{tab:auditoria}
\small
\begin{tabular}{p{0.3\linewidth}p{0.62\linewidth}}
\toprule
Ponto & Decisão e justificativa \\
\midrule
Linguagem do B\&B & C++11 para o núcleo (o Python seria interpretado a cada operação da cota). \\
Estratégia da árvore & DFS em vez de best-first: memória $O(n)$ em vez de fila exponencial. \\
Estado por nó & Incremental (caminho + visitados), sem cópia de matriz por nó. \\
Cota inferior & Medida em pontos aleatórios: com $N=25$ a \textit{naive} estoura 30\,s e a \textit{hk} resolve em 0{,}006\,s; a \textit{hk} só colapsa em $N\approx45$. \\
Isolamento do warm-start & Subgradiente de Held-Karp com passo independente do UB. \\
Timeout & Interno à busca, preservando contadores em NA. \\
MST & Prim denso $O(n^2)$ em vez de Kruskal/heap ($O(n^2\log n)$ em grafo completo). \\
Medição de memória & \texttt{tracemalloc} em passada separada da cronometrada. \\
Memória do B\&B & \texttt{ru\_maxrss} herdado do Python através do \texttt{exec} (35\,MB); remedido isoladamente: 3\,MB. \\
Paralelismo & Processos (não threads, por causa do GIL), um por núcleo físico (\texttt{taskset}), maiores instâncias primeiro, resultados gravados incrementalmente (retomável). \\
Curvas de escala & Medidas até $N=1000$: com $N \le 280$ o overhead fixo por chamada NumPy mascarava o termo assintótico. \\
Feature de agrupamento & Clark-Evans com correção de borda de Donnelly~\cite{clarkevans1954,donnelly1978}. \\
Modelo de sensibilidade & Ridge e validação cruzada por instância (a regressão linear comum sobreajustava: $R^2$ de validação negativo). \\
\bottomrule
\end{tabular}
\end{table}

% ============================================================================
\section{Metodologia Experimental}\label{sec:metodo}

\paragraph{Instâncias.} A TSPLIB~\cite{reinelt1991} não tem instâncias EUC\_2D com
$N \le 40$; as citadas no enunciado são dos tipos GEO, ATT ou EXPLICIT. As
\textit{features} geométricas precisam de coordenadas, então usamos (i) instâncias com coordenadas
(\texttt{NODE\_COORD} ou \texttt{DISPLAY\_DATA}) cujas distâncias foram
\emph{recalculadas} como EUC\_2D, $d(i,j) = \text{nint}(\lVert p_i - p_j \rVert)$, com
coordenadas GEO multiplicadas por 100 para preservar precisão; e (ii) recortes formados
pelos $K$ primeiros pontos de instâncias EUC\_2D maiores, escolhidas por geometria
(uniforme, agrupada, alinhada, em grade). O conjunto tem 15 instâncias, 12 (80\%) com
$N \le 40$ e 3 de estresse (att48, eil51, st70), listadas na Tabela~\ref{tab:instancias}.
O ótimo publicado só vale para eil51 e st70, as únicas originalmente EUC\_2D e sem
recorte; nas demais, o ótimo de referência é o encontrado pelo próprio Branch-and-Bound.
Nas duas, o B\&B chegou aos valores da TSPLIB (426 e 675), o que confere tanto a
implementação quanto o recálculo das distâncias.

\begin{table}[ht]
\centering
\caption{Instâncias de teste. ``Recorte $K$'': os $K$ primeiros pontos da instância
original. O ótimo é o custo EUC\_2D recalculado, provado pelo B\&B com cota \textit{hk}.}
\label{tab:instancias}
\small
\begin{tabular}{lrllr}
\toprule
Instância & $N$ & Origem & Geometria & Ótimo \\
\midrule
kroA100\_k12 & 12 & recorte de kroA100 (EUC\_2D) & uniforme & 9\,775 \\
burma14 & 14 & GEO $\to$ EUC\_2D & geográfica & 3\,087 \\
ulysses16 & 16 & GEO $\to$ EUC\_2D & alongada & 7\,398 \\
berlin52\_k20 & 20 & recorte de berlin52 (EUC\_2D) & urbana & 5\,270 \\
ulysses22 & 22 & GEO $\to$ EUC\_2D & alongada & 7\,529 \\
kroB100\_k25 & 25 & recorte de kroB100 (EUC\_2D) & uniforme & 11\,960 \\
bayg29 & 29 & EXPLICIT (\texttt{DISPLAY\_DATA}) & geográfica & 9\,073 \\
bier127\_k30 & 30 & recorte de bier127 (EUC\_2D) & agrupada & 17\,900 \\
a280\_k33 & 33 & recorte de a280 (EUC\_2D) & alinhada (circuito) & 457 \\
pr76\_k35 & 35 & recorte de pr76 (EUC\_2D) & irregular & 44\,237 \\
lin105\_k38 & 38 & recorte de lin105 (EUC\_2D) & grade & 5\,568 \\
eil101\_k40 & 40 & recorte de eil101 (EUC\_2D) & uniforme & 413 \\
att48 & 48 & ATT $\to$ EUC\_2D & geográfica & 33\,522 \\
eil51 & 51 & EUC\_2D & uniforme & 426 \\
st70 & 70 & EUC\_2D & uniforme & 675 \\
\bottomrule
\end{tabular}
\end{table}

\paragraph{Ambiente.} Intel Core i5-10210U (4 núcleos físicos), 7\,GB de RAM, WSL2/Ubuntu
22.04, g++ 11.4 (\texttt{-O3 -march=native}), Python 3.10, NumPy 2.2, NetworkX 3.4,
scikit-learn 1.7. Três execuções do B\&B em paralelo, cada uma fixada em um núcleo físico
distinto (\texttt{taskset}, CPUs lógicas 0, 2 e 4), com prioridade reduzida
(\texttt{nice -n 10}). O quarto núcleo ficou livre porque a máquina continuou em uso
durante a rodada. Com um processo por núcleo físico, dois processos nunca dividem o
mesmo núcleo por \textit{hyper-threading}, o que alteraria os tempos. As 150 execuções
levaram 5\,h\,21\,min.

\paragraph{Protocolo.} Para cada instância, estratégia de \textit{warm-start} e cota
(\textit{hk} e \textit{mst}), uma execução do B\&B com limite de 30\,min. O tempo total
de cada estratégia é o tempo da heurística mais o tempo da busca. Para as curvas de
complexidade das heurísticas, usamos recortes de $N=50$ a $1000$ da pr1002, com passo
50 e mediana de 3 repetições. A cota \textit{naive} só foi usada nos testes preliminares
com pontos aleatórios (Tabela~\ref{tab:auditoria}): como ela colapsa por volta de
$N \approx 25$, ficou fora da rodada principal.

\paragraph{Ajustes de complexidade.} Para cada sub-rotina, ajustamos
$t = a + b\,g(N)$ com $g \in \{N, N\log N, N^2, N^2\log N, N^3\}$ por mínimos quadrados
não negativos. Cada observação tem peso $1/t$, de modo que o ajuste minimiza o erro
relativo; sem esse peso, os poucos pontos com $N$ grande decidiriam o ajuste sozinhos.
Escolhemos o modelo de maior $R^2$ em escala log. Ajustamos também a lei de potência
$t = aN^b$ por regressão em log-log, com intervalo de confiança de 95\% para $b$.

\paragraph{Sensibilidade.} O gap de cada aproximação é $(\text{custo} -
\text{ótimo})/\text{ótimo}$. Calculamos a correlação de Spearman de cada
\textit{feature} com o gap de cada algoritmo. Depois treinamos uma árvore de decisão
(profundidade 3) e uma regressão Ridge, validadas deixando uma instância inteira de
fora, em duas versões: \emph{bruta}, com o algoritmo como variável indicadora, e
\emph{centrada}, em que se subtrai de cada gap a média do seu algoritmo e o modelo usa só
as \textit{features}.

% ============================================================================
\section{Resultados}\label{sec:res}

\subsection{Complexidade empírica}

\begin{figure}[ht]
\centering
\includegraphics[width=\linewidth]{complexidade_tempo}
\caption{Tempo das sub-rotinas (esq.) e dos algoritmos completos (dir.) em função de $N$,
com a curva teórica de melhor ajuste.}
\label{fig:complexidade}
\end{figure}

\begin{table}[ht]
\centering
\caption{Melhor ajuste por sub-rotina e expoente da lei de potência $t = a N^b$.}
\label{tab:complexidade}
\restable{complexidade_melhor_ajuste}
\end{table}

A Tabela~\ref{tab:complexidade} compara cada sub-rotina com a sua complexidade teórica.
As duas sub-rotinas cúbicas se comportam como a teoria prevê. O emparelhamento tem
expoente $b = 2{,}70 \pm 0{,}09$ e melhor ajuste $N^3$ ($R^2_{\log} = 0{,}985$), como
se espera do blossom $O(k^3)$ sobre os $k$ vértices ímpares da MST, já que $k$ cresce
proporcionalmente a $N$. O expoente fica abaixo de 3 porque termos de ordem menor ainda
pesam na faixa medida. A inserção mais barata tem $b = 2{,}50 \pm 0{,}22$ e melhor
ajuste $N^3$, e a construção do grafo do emparelhamento, quadrática em $k$, tem
$b = 2{,}01$.

Já as sub-rotinas quadráticas vetorizadas ficam longe do expoente teórico até
$N = 1000$: o Prim denso tem $b = 1{,}27$ e o vizinho mais próximo, $b = 1{,}32$,
ambos $O(N^2)$ na teoria. Cada uma faz $N$ iterações, e cada iteração chama o NumPy
sobre um vetor de tamanho $N$, com custo $c_0 + c_1 N$. Nessa faixa, o
\textit{overhead} fixo $c_0$ de cada chamada (alguns microssegundos) ainda é maior que
o trabalho vetorial $c_1 N$. É por isso que o critério de seleção escolhe $N \log N$
para o Prim; com $N$ maior, o termo quadrático apareceria. O ruído na escala de
microssegundos também explica o $R^2$ linear baixo do vizinho mais próximo (0,52). Em
números absolutos, o Twice-around-the-tree completo leva 0,03\,s em $N = 1000$, e o
Christofides, 66\,s.

\begin{figure}[ht]
\centering
\includegraphics[width=0.75\linewidth]{complexidade_memoria}
\caption{Pico de memória (\texttt{tracemalloc}) das sub-rotinas em função de $N$.}
\label{fig:memoria}
\end{figure}

\begin{table}[ht]
\centering
\caption{Lei de potência do pico de memória ($m = aN^b$) e pico em $N=1000$.}
\label{tab:memoria}
\restable{complexidade_memoria}
\end{table}

Na memória (Figura~\ref{fig:memoria}, Tabela~\ref{tab:memoria}), a diferença entre as
representações é ainda maior que no tempo. A MST guarda só vetores de tamanho $N$
($b = 1{,}17$, 68\,KB em $N=1000$). A inserção mantém uma matriz $N \times N$ de custos
de inserção ($b = 1{,}96$, 10\,MB). O grafo do emparelhamento no NetworkX cresce com
expoente parecido ($b = 2{,}06$), mas ocupa 39\,MB, porque cada uma das
$\approx k^2/2$ arestas é um dicionário Python e custa muito mais que uma entrada de
matriz.

\begin{figure}[ht]
\centering
\includegraphics[width=\linewidth]{bnb_crescimento}
\caption{Nós expandidos (esq.) e tempo (dir.) do B\&B com \textit{warm-start}
Christofides, em escala log, com o ajuste exponencial de cada cota. Marcadores vazados:
execuções interrompidas em 30\,min, fora do ajuste.}
\label{fig:crescimento}
\end{figure}

\begin{table}[ht]
\centering
\caption{Crescimento exponencial do B\&B (\textit{warm-start} Christofides): fator
multiplicativo por vértice adicional, ajustado sobre as instâncias resolvidas.}
\label{tab:crescimento}
\restable{bnb_crescimento}
\end{table}

O tempo do Branch-and-Bound cresce exponencialmente com $N$
(Figura~\ref{fig:crescimento}, Tabela~\ref{tab:crescimento}): cada vértice a mais
multiplica o tempo por 1,28, em média, com a cota \textit{hk}, e por 1,40 com a
\textit{mst}. Os pontos se espalham bastante em torno da reta (RMSE de 0,55 décadas
para a \textit{hk} e de 1,0 para a \textit{mst}), porque a geometria pesa tanto quanto
$N$. Com a cota \textit{hk}, a lin105\_k38 ($N=38$) é resolvida em 0,03\,s e a
eil101\_k40 ($N=40$) em 1,5\,s; a eil51 leva 5\,s e a st70, 609\,s. Extrapolando o
ajuste a partir da st70, a cota \textit{hk} chegaria ao limite de 30\,min por volta de
$N \approx 75$. Com a \textit{mst}, o colapso acontece dentro do próprio conjunto de
teste e varia com a instância: a ulysses22 ($N=22$) não termina em 30\,min, e a
lin105\_k38 ($N=38$) termina em 6\,min.

\subsection{Impacto do warm-start na árvore de busca}

\begin{figure}[ht]
\centering
\includegraphics[width=\linewidth]{warmstart_nos}
\caption{Nós expandidos por instância e estratégia de \textit{warm-start}, para as cotas
\textit{hk} e \textit{mst}. Barras hachuradas: execução interrompida em 30\,min (NA).}
\label{fig:warmstart}
\end{figure}

\begin{table}[ht]
\centering
\caption{Resumo por cota e estratégia (médias geométricas sobre instâncias resolvidas por todas as estratégias).}
\label{tab:warmstart}
\restable{warmstart_resumo}
\end{table}

\begin{figure}[ht]
\centering
\includegraphics[width=0.6\linewidth]{warmstart_tradeoff}
\caption{Gap do limite superior inicial contra a razão de nós expandidos em relação ao
início trivial (cota \textit{hk}, eixo $x$ em escala symlog).}
\label{fig:tradeoff}
\end{figure}

\paragraph{Vizinho mais próximo e Twice-around-the-tree.} Com a cota \textit{hk}, essas
duas estratégias expandem 99,5\% e 99,8\% dos nós do início trivial
(Tabela~\ref{tab:warmstart}), embora partam de tours com gap médio de 21\% e 28\%,
contra 146\% do trivial. Como a DFS visita os filhos do mais próximo para o mais
distante, o primeiro ramo que ela percorre é o tour do vizinho mais próximo a partir do
vértice 0. Depois de $N$ nós, a busca trivial já tem esse limite superior, e o
\textit{warm-start} NN economiza só essa primeira descida; na bier127\_k30, são 8 nós a
menos em 306\,303. O Twice-around-the-tree é pior que o vizinho mais próximo em 12 das
15 instâncias. Nesses casos, a primeira descida já encontra um tour melhor que o dele,
e em 11 dos 12 a árvore fica igual à do trivial (306\,303 nós também na bier127\_k30);
na a280\_k33, a diferença é de 6 nós em 533.

\paragraph{Fecho convexo.} A inserção a partir do fecho convexo produz tours com gap
médio de 2,2\% e acerta o ótimo em 4 das 15 instâncias. Nessas 4, só a raiz é
expandida, porque todos os filhos já têm $\lceil LB \rceil \ge UB$. Em média
geométrica, a estratégia expande 8,4\% dos nós do trivial e reduz o tempo total
(heurística mais busca) em 2,1 vezes. Na st70, a instância mais difícil, a árvore cai
de 3,8 milhões para 151 mil nós, e o tempo, de 541\,s para 69\,s. Na
Figura~\ref{fig:tradeoff}, a poda só melhora com gap inicial abaixo de cerca de 10\%,
que é aproximadamente o gap do tour que a DFS encontra sozinha.

\paragraph{Christofides.} Com gap médio de 8,7\%, o Christofides reduz os nós para 67\%
dos do trivial, e em alguns casos a redução é grande: 24 vezes na bier127\_k30
(306\,303 $\to$ 12\,580 nós). No tempo total, o ganho é de só 3\%. A cota \textit{hk}
resolve a maior parte das instâncias em milissegundos, e os 6,6\,ms do emparelhamento
consomem o que a busca economiza. Esse \textit{warm-start} só compensa quando a busca
demora o bastante para amortizar a heurística. A inserção leva 1,5\,ms e dá o melhor
limite, então ganha nos dois critérios.

\paragraph{Cota mst.} Com a cota \textit{mst}, a estratégia de \textit{warm-start} quase
não faz diferença. O fecho convexo, a melhor delas, ainda expande 83\% dos nós do
trivial, e as 5 instâncias que esgotam o tempo esgotam com todas as 5 estratégias. Na
raiz, a \textit{hk} fica de 0,1\% a 4,9\% abaixo do ótimo, e a \textit{mst}, de 10\%
a 33\%. Com uma cota tão frouxa, a maioria dos nós tem $LB$ bem abaixo do ótimo, e
nenhum limite superior os poda, nem o próprio ótimo. O caso extremo é a ulysses22, que
tem a maior folga da \textit{mst} (33\%) e a menor da \textit{hk} (0,1\%). A
\textit{hk} a resolve em 8\,ms; a \textit{mst} expande $8{,}4 \times 10^8$ nós em
30\,min sem provar a otimalidade, embora já tivesse o ótimo (7\,529) como incumbente.
Nas outras execuções interrompidas, o incumbente final varia: a eil51 termina a 1
unidade do ótimo, e a st70 termina 8,7\% acima com o início trivial e 3,9\% acima com o
fecho convexo.

\subsection{Análise de sensibilidade}

As \textit{features} descrevem a geometria de cada instância. \texttt{aspect\_ratio} é a
razão entre os lados da caixa envolvente. \texttt{hull\_ratio} é a fração de pontos no
fecho convexo. \texttt{clark\_evans\_R} é o índice de agrupamento de Clark-Evans:
abaixo de 1 indica pontos agrupados, acima de 1, pontos regulares. A janela usada é o
fecho, o que dá ao índice um viés positivo para $N$ pequeno; por isso, ele só serve para
comparar instâncias.
\texttt{spatial\_entropy} é a entropia normalizada da ocupação de uma grade.
\texttt{nn\_cv} e \texttt{dist\_cv} são os coeficientes de variação das distâncias ao
vizinho mais próximo e de todas as distâncias. As três \textit{features}
\texttt{mst\_*} vêm da MST: o grau máximo, a fração de vértices de grau ímpar e o
coeficiente de variação das arestas.

\begin{table}[ht]
\centering
\caption{Gap de cada algoritmo em relação ao ótimo ((custo $-$ ótimo)/ótimo).}
\label{tab:gaps}
\restable{gaps}
\end{table}

Os gaps ficaram bem abaixo das garantias teóricas (Tabela~\ref{tab:gaps}). O pior caso do
Twice-around-the-tree é 42\%, na eil51, para uma garantia de 100\%; o do Christofides é
18\%, na eil101\_k40, para uma garantia de 50\%. Em média, o Christofides fica a
8,7\% do ótimo e o Twice-around-the-tree a 27,7\%. A heurística de inserção a partir do
fecho convexo, sem garantia constante, é melhor que as duas em 13 das 15 instâncias,
com média de 2,2\%. As exceções são ulysses16 e lin105\_k38, em que o Christofides a
supera por menos de 0,4 ponto percentual.

\begin{figure}[ht]
\centering
\includegraphics[width=0.75\linewidth]{sensibilidade_spearman}
\caption{Correlação de Spearman entre \textit{features} e gap de aproximação (* $p<0{,}05$).}
\label{fig:spearman}
\end{figure}

\begin{table}[ht]
\centering
\caption{Desempenho dos modelos de sensibilidade (validação cruzada deixando uma instância de fora).}
\label{tab:modelos}
\restable{sensibilidade_modelos}
\end{table}

\paragraph{Correlações.} Na Figura~\ref{fig:spearman}, cada algoritmo responde a
\textit{features} diferentes. O gap do Christofides é maior em instâncias espalhadas de
modo uniforme: tem correlação $+0{,}83$ com a entropia espacial e $-0{,}87$ com a
dispersão das distâncias (\texttt{dist\_cv}), e cresce com a fração de vértices ímpares
na MST ($+0{,}63$). Com mais vértices ímpares, o emparelhamento acrescenta mais arestas
ao circuito euleriano, e o circuito precisa de mais atalhos. Em instâncias agrupadas ou
alongadas, a MST se parece mais com um caminho e tem menos vértices ímpares. O gap do
Twice-around-the-tree e o da inserção a partir do fecho crescem com $N$ ($+0{,}70$ e
$+0{,}71$) e caem em instâncias alongadas (\texttt{aspect\_ratio}, de $-0{,}60$ a
$-0{,}64$, o mesmo valendo para o Christofides). Quando os pontos estão quase
alinhados, a pré-ordem da MST se aproxima de um percurso de ida e volta, que é
quase ótimo. O vizinho mais próximo não tem nenhuma correlação significativa.

\paragraph{Casos específicos.} A a280\_k33, uma placa de circuito com pontos alinhados
e a maior razão de aspecto (4,95), tem os resultados mais desiguais: o Christofides
fica a 2,4\% e o fecho a 0\%, enquanto o vizinho mais próximo e o
Twice-around-the-tree ficam a 31\% e 32\%. Com os pontos em fileiras, as escolhas
gulosas deixam vértices isolados para trás. A lin105\_k38, em grade, tem muitas
distâncias empatadas, mas o Christofides fica a 3,6\% e o B\&B com cota \textit{hk} a
resolve em 0,03\,s. O pior caso do Christofides (18\%) é a eil101\_k40, uniforme e com
entropia de 0,98.

\paragraph{Modelos preditivos.} Na versão bruta, árvore e Ridge chegam a $R^2 \approx
0{,}54$ na validação cruzada (Tabela~\ref{tab:modelos}), mas quase todo esse poder vem
da identidade do algoritmo: na árvore, as indicadoras de Twice-around-the-tree e de
vizinho mais próximo somam 77\% da importância. A versão centrada testa se a geometria
explica o desvio de cada algoritmo em relação à sua própria média, e o $R^2$ de
validação é $\approx -0{,}03$, sem poder preditivo fora da amostra. A árvore ainda
separa as instâncias por $N \le 23$ e usa \texttt{mst\_edge\_cv} e
\texttt{mst\_odd\_frac}, mas esse ajuste ($R^2 = 0{,}29$ no treino) não generaliza.

Com 15 instâncias, esse resultado já era provável, e as correlações acima servem como
hipóteses a testar. As \textit{features} são correlacionadas entre si
(\texttt{hull\_ratio}, por exemplo, cai naturalmente com $N$), e cada algoritmo responde
a \textit{features} diferentes, o que enfraquece um modelo único. Além disso, foram 40
testes de Spearman, e ao nível de 5\% cerca de 2 dariam significativos só por acaso.
Para confirmar os efeitos, seria preciso gerar instâncias com geometria controlada,
variando uma \textit{feature} por vez.

% ============================================================================
\section{Conclusões}\label{sec:concl}

Com as penalidades de Held-Karp calculadas uma única vez na raiz, o B\&B resolveu as 15
instâncias, até $N=70$, dentro do limite de 30\,min. Com a cota de MST pura, 5
instâncias esgotaram o tempo, a primeira com $N=22$. Na raiz, a folga da \textit{hk}
foi de 0,1\% a 4,9\%, e a da \textit{mst}, de 10\% a 33\%; essa diferença pesou mais
na poda que qualquer \textit{warm-start}.

Numa DFS que visita primeiro o vizinho mais próximo, iniciar com o tour do vizinho mais
próximo ou do Twice-around-the-tree não muda a busca, porque ela encontra um tour pelo
menos tão bom na primeira descida. Só limites melhores que esse tour guloso aumentaram a
poda. A inserção a partir do fecho convexo, barata e quase ótima, reduziu a árvore em
12 vezes. O Christofides reduziu os nós em 33\%, mas a busca nessas instâncias é tão
rápida que o custo do emparelhamento anulou o ganho de tempo.

O emparelhamento ocupa mais de 94\% do tempo do Christofides e quase toda a sua memória,
com crescimento empírico $N^{2,7}$. As sub-rotinas quadráticas vetorizadas ainda não
mostram o expoente teórico até $N=1000$, por causa do \textit{overhead} de cada chamada
ao NumPy.

Na prática, os gaps ficaram bem abaixo das garantias: 18\% no pior caso do Christofides
e 42\% no do Twice-around-the-tree. As correlações entre geometria e gap fazem sentido
para cada algoritmo, mas 15 instâncias não bastaram para um modelo preditivo.

Três extensões seguem destes resultados: atualizar as penalidades de Held-Karp ao longo
da árvore, para levar o colapso além de $N \approx 75$; trocar o emparelhamento do
NetworkX por uma implementação compilada ou gulosa; e gerar instâncias sintéticas com
geometria controlada para a análise de sensibilidade.

\bibliographystyle{sbc}
\bibliography{refs}

\end{document}
